
\begin{derivation}[从\eqnrefs{eq:phi4-si-coupling-r10}到\eqnrefs{eq:si-reduced-vertices-r10}]

正文\eqnrefs{eq:phi4-si-coupling-r10} 中显式写出的 $1/4!$ 不是图规则的附加约定，而是为了抵消同一个局域顶角上四个相同标量场的等价 Wick 配对数。逐个计数这 $4!$ 种收缩并截去外腿后，就得到正文\eqnrefs{eq:si-reduced-vertices-r10} 中的 $-\ii g_4$。



SI 相互作用密度已经固定为
\begin{equation*}
 \Lag_I(z)=-\frac{g_4}{4!\,\hbar c}\varphi^4(z).
\end{equation*}
因此 Dyson 一阶对四点 Green 函数的贡献与
\begin{equation*}
 \frac{1}{4!}
 \int\dd^4 z\,
 \langle0|\Tord\{
 \varphi(x_1)\varphi(x_2)\varphi(x_3)\varphi(x_4)
 \varphi^4(z)
 \}|0\rangle
 \end{equation*}
成正比；所有 $\hbar,c$ 的共同归一化已在
\eqnrefs{eq:si-reduced-vertices-r10}
的约化过程中统一提出。要得到连通树图，每个外部场必须与 $z$ 点四个场中的一个收缩。第一条外部场有 4 个选择，第二条剩 3 个，第三条剩 2 个，最后剩 1 个，所以同一个收缩乘积
\begin{equation*}
 \Delta_{1z}\Delta_{2z}\Delta_{3z}\Delta_{4z}
\end{equation*}
出现
\begin{equation*}
 4\times3\times2\times1=4!
\end{equation*}
次，恰好抵消拉氏量中的 $1/4!$。因此截肢外腿并提出统一 SI 归一化后，约化顶角只剩
\begin{equation*}
 {-\ii g_4.}
\end{equation*}
$4!$ 不是被经验地“删掉”，而是被等价 Wick 收缩精确抵消。
\end{derivation}

\begin{derivation}[从\eqnrefs{eq:qed-gell-mann-low-dp7}到\eqnrefs{eq:phi4-bubble-symmetry-factor-dp68}]
考虑 $\varphi^4$ 理论四点函数的二阶 Dyson 项
\begin{equation*}
 \frac{1}{2!}
 \left(-\frac{\ii g_4}{4!\,\hbar^2c^2}\right)^2
 \int\dd^4z\dd^4w\,
 \langle0|\mathcal T\{
 \varphi_1\varphi_2\varphi_3\varphi_4
 \varphi^4(z)\varphi^4(w)\}|0\rangle .
\end{equation*}
只数某一个固定 $s$ 道拓扑：$x_1,x_2$ 接到一个顶角，$x_3,x_4$ 接到另一个顶角，两个顶角之间有两条内部线。不能只说“交换内部线给 $2!$”，因为还要把 Dyson 的 $2!$ 和每个顶角的 $4!$ 一起计数。

先决定哪一个带积分标签的顶角接 $(x_1,x_2)$：有 $2$ 种选择。固定以后，在该顶角的四个同点场中，为有标签的 $x_1$ 选收缩对象有 $4$ 种，为 $x_2$ 再选有 $3$ 种，所以有 $4\times3=12$ 种。另一顶角给 $(x_3,x_4)$ 同样有 $12$ 种。每个顶角此时都剩两个场，两条跨顶角内部线的配对还有 $2!$ 种。因此产生同一个带标签 $s$ 道图的 Wick 收缩总数为
\begin{equation*}
 N_s=2\times12\times12\times2=576.
\end{equation*}
而 Dyson 与两个相互作用顶角的组合分母是
\begin{equation*}
 2!(4!)^2=2\times24^2=1152.
\end{equation*}
所以净组合系数严格为
\begin{equation*}
 \frac{N_s}{2!(4!)^2}=\frac{576}{1152}=\frac12.
\end{equation*}
这就是所谓泡图对称因子。它不是在画完图以后经验性补上的数，而是完整 Dyson--Wick 展开中“产生同一带标签图的收缩数”除以“指数展开和拉氏量中的排列分母”所得的商。

截肢四条外腿后，$s$ 道一圈振幅因而为
\begin{equation*}
 \ii\mathcal M_s^{(1)}
 =\frac12(-\ii g_4)^2
 \int_k^{(4)}
 \frac{\ii}{k^2-m^2c^2+\ii0^+}
 \frac{\ii}{(P-k)^2-m^2c^2+\ii0^+},
 \qquad P=p_1+p_2,
\end{equation*}
即正文\eqnrefs{eq:phi4-bubble-symmetry-factor-dp68}。$t$、$u$ 道来自外部标签的另外两种不等价分组，而不是这个 $1/2$ 内部对称因子的额外来源。
\end{derivation}

\begin{derivation}[从\eqnrefs{eq:qed-gell-mann-low-dp7}到\eqnrefs{eq:identical-fermion-exchange-sign-dp68}]
取 SI Yukawa 相互作用

\begin{equation*}
 \Lag_I=-y\sqrt{\hbar c}\,\varphi\bar\psi\psi,
\end{equation*}
考虑
\begin{equation*}
 f(p_1,s_1)+f(p_2,s_2)
 \to
 f(p_3,s_3)+f(p_4,s_4).
\end{equation*}
固定两电子初末态的算符次序：
\begin{align*}
 |i\rangle
 &=b_{s_1}^\dagger(p_1)b_{s_2}^\dagger(p_2)|0\rangle,\\
 \langle f|
 &=\langle0|b_{s_4}(p_4)b_{s_3}(p_3).
\end{align*}
二阶 Dyson 项中两个媒介粒子场收缩成内部标量线，每个 $\bar\psi\psi$ 顶角各提供一只电子湮灭算符和一只产生算符。第一种配对保持 $1\to3$、$2\to4$，产生
\begin{equation*}
 [\bar u_3u_1][\bar u_4u_2]
 \Delta_{\rm red}(p_1-p_3).
\end{equation*}
第二种配对交换末态标签：$1\to4$、$2\to3$。要把同一外部-态次序重排成两条标准旋量 chains，必须额外交换一次费米算符；由
\begin{equation*}
 b_a b_b^\dagger
 =-b_b^\dagger b_a+\{b_a,b_b^\dagger\}
\end{equation*}
得到额外负号，故第二项为
\begin{equation*}
 -[\bar u_4u_1][\bar u_3u_2]
 \Delta_{\rm red}(p_1-p_4).
\end{equation*}
于是
\begin{equation*}
 {\mathcal M=\mathcal M_t-\mathcal M_u.}
\end{equation*}
这个负号完全来自 Fock 空间反对易关系，与是否使用 SI 或约化图因子无关。
\end{derivation}

\begin{derivation}[从\eqnrefs{eq:spin-sum-si-r9}到\eqnrefs{eq:qed-basic-tensor-si-r9}]

未偏振截面需要对末态自旋求和、对初态自旋平均。正文\eqnrefs{eq:spin-sum-si-r9} 把这些离散自旋指标变成 $\slashed p+mc$，随后利用迹循环性可把两条外部自旋链合并成单个 Dirac 迹；对矢量流执行这一步便得到正文\eqnrefs{eq:qed-basic-tensor-si-r9}。



取
\begin{equation*}
 B_{rs}=\bar u_r(p')\mathcal O u_s(p),
\end{equation*}
其中 $\mathcal O$ 是任意 $4\times4$ 旋量空间矩阵；这里的花体 $\mathcal O$ 只表示插在两个外线自旋子之间的算符，不与 1PI 顶角 $\Gamma^\mu$ 或衰变率记号混用。先求复共轭：
\begin{align*}
 B_{rs}^*
 &=[u_r^\dagger(p')\gamma^0\mathcal O u_s(p)]^*,\\
 &=u_s^\dagger(p)\mathcal O^\dagger\gamma^0u_r(p'),\\
 &=\bar u_s(p)\gamma^0\mathcal O^\dagger\gamma^0u_r(p').
\end{align*}
定义
\begin{equation*}
 \overline{\mathcal O}\equiv\gamma^0\mathcal O^\dagger\gamma^0.
\end{equation*}
于是
\begin{align*}
 \sum_{r,s}|B_{rs}|^2
 &=\sum_{r,s}
 \bar u_r(p')\mathcal O u_s(p)
 \bar u_s(p)\overline{\mathcal O} u_r(p'),\\
 &=\sum_r\bar u_r(p')\mathcal O
 \left[\sum_su_s(p)\bar u_s(p)\right]
 \overline{\mathcal O} u_r(p'),\\
 &=\sum_r\bar u_r(p')\mathcal O(\slashed p+mc)
 \overline{\mathcal O} u_r(p').
\end{align*}
把最后一个标量写成迹：
\begin{equation*}
 \bar u_r X u_r=\Tr[Xu_r\bar u_r],
\end{equation*}
再对 $r$ 求和，
\begin{align*}
 \sum_{r,s}|B_{rs}|^2
 &=\Tr\!\left[
 \mathcal O(\slashed p+mc)\overline{\mathcal O}
 \sum_ru_r(p')\bar u_r(p')
 \right],\\
 &={
 \Tr\!\left[(\slashed p'+mc)\mathcal O
 (\slashed p+mc)\overline{\mathcal O}\right].}
 \end{align*}
未偏振入射态由自旋密度矩阵 $\mathbb I_2/2$ 表示，末态自旋求和由完备关系实现；两者把有限个外线自旋指标精确收缩成矩阵迹。

对矢量流取 $\mathcal O=\gamma^\mu$，并使用
\begin{equation*}
 \Tr(\gamma^\alpha\gamma^\beta)=4\eta^{\alpha\beta},
 \qquad
 \Tr(\gamma^\alpha\gamma^\mu\gamma^\beta\gamma^\nu)
 =4(\eta^{\alpha\mu}\eta^{\beta\nu}
 -\eta^{\alpha\beta}\eta^{\mu\nu}
 +\eta^{\alpha\nu}\eta^{\mu\beta}),
\end{equation*}
奇数个 $\gamma$ 矩阵的迹为零，于是
\begin{align*}
 \Tr[(\slashed a+mc)\gamma^\mu(\slashed b+mc)\gamma^\nu]
 &=\Tr(\slashed a\gamma^\mu\slashed b\gamma^\nu)
 +m^2c^2\Tr(\gamma^\mu\gamma^\nu),\\
 &=4[a^\mu b^\nu+a^\nu b^\mu-(a\!\cdot b-m^2c^2)\eta^{\mu\nu}],
\end{align*}
即\eqnrefs{eq:qed-basic-tensor-si-r9}。
\end{derivation}

\begin{derivation}[从\eqnrefs{eq:spin-fourvector-r9}到\eqnrefs{eq:polarized-spinor-projector-r9}]
Landau--Lifshitz 型的偏振散射计算经常不先做自旋求和，而是保留一个指定的极化 four-矢量。先在粒子静止系取单位三矢量 $\boldsymbol\zeta$ 指定自旋方向，定义

\begin{equation*}
 s_{\rm rest}^\mu=(0,\boldsymbol\zeta),
 \qquad
 s_{\rm rest}^2=-1,
 \qquad
 s_{\rm rest}\!\cdot p_{\rm rest}=0.
\end{equation*}
把 $p_{\rm rest}^\mu=(mc,\mathbf0)$ 与 $s_{\rm rest}^\mu$ 用同一个 Lorentz 推动变换到一般参考系。若 $\mathbf v$ 是粒子速度、$\gamma=E/(mc^2)$，推动给
\begin{equation*}
 {
 s^0=\gamma\frac{\mathbf v\!\cdot\boldsymbol\zeta}{c},
 \qquad
 \mathbf s
 =\boldsymbol\zeta
 +\frac{\gamma^2}{\gamma+1}
 \frac{\mathbf v(\mathbf v\!\cdot\boldsymbol\zeta)}{c^2}.}
 \end{equation*}
直接代入可检查
\begin{equation*}
 s^2=-1,
 \qquad
 s\!\cdot p=0.
\end{equation*}
在静止系，正能旋量的自旋投影算符是
\begin{equation*}
 \frac12(1+\boldsymbol\Sigma\!\cdot\boldsymbol\zeta).
\end{equation*}
利用
\begin{equation*}
 \gamma^5\slashed s_{\rm rest}
 =\boldsymbol\Sigma\!\cdot\boldsymbol\zeta
\end{equation*}
和正能投影算符 $\slashed p_{\rm rest}+mc=mc(\gamma^0+1)$，单一自旋态满足
\begin{equation*}
 u(p_{\rm rest},s)\bar u(p_{\rm rest},s)
 =\frac12(\slashed p_{\rm rest}+mc)
 (1+\gamma^5\slashed s_{\rm rest}).
\end{equation*}
两边在 Lorentz 变换下都按 $S(\Lambda)(\cdots)S^{-1}(\Lambda)$ 变换，所以一般参考系中
\begin{equation*}
 {
 u(p,s)\bar u(p,s)
 =\frac12(\slashed p+mc)(1+\gamma^5\slashed s).}
 \end{equation*}
把相反偏振 $s^\mu\to-s^\mu$ 的两式相加，$\gamma^5\slashed s$ 项消失，所得和等于正文的未偏振完备关系。
\end{derivation}

\begin{derivation}[从\eqnrefs{eq:spin-fourvector-r9}到\eqnrefs{eq:electron-helicity-projector-dp9}]

螺旋度把自旋量子化轴选在动量方向，因此仍属于同一自旋自由度的动量适配基。将正文\eqnrefs{eq:spin-fourvector-r9} 中的静止系自旋方向取为 $\boldsymbol\zeta=\sigma_h\widehat{\mathbf p}$，先得到对应的协变自旋四矢量，再代入固定自旋密度矩阵，就得到正文\eqnrefs{eq:electron-helicity-projector-dp9}。



令 $\sigma_h=\pm1$ 表示两倍螺旋度符号，即电子螺旋度本征值为
\begin{equation*}
 h=\frac{\sigma_h}{2}\hbar.
\end{equation*}
对 $\mathbf p\neq0$ 定义 $\widehat{\mathbf p}=\mathbf p/|\mathbf p|$。把静止系自旋方向取为 $\boldsymbol\zeta=\sigma_h\widehat{\mathbf p}$，代入\eqnrefs{eq:spin-fourvector-r9} 的推动公式 \eqnrefs{eq:spin-fourvector-r9}。因为推动与自旋方向共线，得到
\begin{equation*}
 {
 s_h^\mu
 =\sigma_h\left(
 \frac{|\mathbf p|}{m_{\mathrm e} c},
 \frac{E_{\mathbf p}}{m_{\mathrm e} c^2}\widehat{\mathbf p}
 \right).}
 \end{equation*}
逐项检查：
\begin{align*}
 p\!\cdot s_h
 &=\frac{E_{\mathbf p}}{c}\frac{\sigma_h|\mathbf p|}{m_{\mathrm e} c}
 -\mathbf p\cdot
 \frac{\sigma_hE_{\mathbf p}}{m_{\mathrm e} c^2}\widehat{\mathbf p}
 =0,\\
 s_h^2
 &=\frac{|\mathbf p|^2}{m_{\mathrm e}^2c^2}
 -\frac{E_{\mathbf p}^2}{m_{\mathrm e}^2c^4}
 =-1,
\end{align*}
其中最后一步用了 $E_{\mathbf p}^2=m_{\mathrm e}^2c^4+c^2|\mathbf p|^2$。所以它确实是 \eqnrefs{eq:polarized-spinor-projector-r9} 所要求的单位类空自旋四矢量。

直接代入偏振密度矩阵，得到螺旋度可分辨的外部电子投影算符
\begin{equation*}
 {
 u_h(p)\bar u_h(p)
 =\frac12(\slashed p+m_{\mathrm e} c)
 \left(1+\gamma^5\slashed s_h\right).}
 \end{equation*}
把 $\sigma_h=+1$ 与 $-1$ 相加，$\slashed s_h$ 项相消，恢复
$\sum_hu_h\bar u_h=\slashed p+m_{\mathrm e} c$。因此未偏振自旋求和就是两个螺旋度投影算符的完备和。

在 $\mathbf p\to0$ 时 $\widehat{\mathbf p}$ 无定义，这不是公式缺陷：有质量粒子的静止系没有优选传播方向，螺旋度本来就不是静止系量子数。此时应回到任意静止系自旋轴的 $SU(2)$ 基。
\end{derivation}

\begin{derivation}[从\eqnrefs{eq:physical-photon-polarization-projector-dp42}到\eqnrefs{eq:photon-helicity-projector-covariant-dp9}]
取零质量光子动量 $k^\mu$，并选任意满足 $k\!\cdot n\neq0$ 的辅助四矢量 $n^\mu$。由\eqnrefs{eq:physical-photon-polarization-projector-dp42} 得到物理横向偏振求和

\begin{equation*}
 \mathsf P_\perp^{\mu\nu}(k;n)
 \equiv
 -\eta^{\mu\nu}
 +\frac{k^\mu n^\nu+n^\mu k^\nu}{k\!\cdot n}
 -\frac{n^2k^\mu k^\nu}{(k\!\cdot n)^2}.
 \tag{\ref{eq:physical-photon-polarization-projector-dp42}}
\end{equation*}
在该二维物理子空间选两个实正交偏振矢量 $e_1^\mu,e_2^\mu$，使
\begin{equation*}
 e_a\!\cdot k=e_a\!\cdot n=0,
 \qquad
 e_a\!\cdot e_b=-\delta_{ab},
 \qquad
 \mathsf P_\perp^{\mu\nu}=e_1^\mu e_1^\nu+e_2^\mu e_2^\nu.
\end{equation*}
定义圆偏振（螺旋度）基
\begin{equation*}
 \varepsilon_\lambda^\mu(k;n)
 \equiv\frac1{\sqrt2}
 \left(e_1^\mu+\ii\lambda e_2^\mu\right),
 \qquad \lambda=\pm1.
 \end{equation*}
于是
\begin{align*}
 \varepsilon_\lambda^\mu\varepsilon_\lambda^{\nu *}
 &=\frac12\Bigl[
 e_1^\mu e_1^\nu+e_2^\mu e_2^\nu
 -\ii\lambda(e_1^\mu e_2^\nu-e_2^\mu e_1^\nu)
 \Bigr].
 \end{align*}
取空间取向与 $\epsilon^{0123}=+1$ 一致，二维面积形式可协变化为
\begin{equation*}
 \mathsf C^{\mu\nu}(k;n)
 \equiv
 \frac{\epsilon^{\mu\nu\rho\sigma}k_\rho n_\sigma}
 {k\!\cdot n}.
 \end{equation*}
因此每个光子螺旋度的秩一投影算符是
\begin{equation*}
 {
 \mathsf P_\lambda^{\mu\nu}(k;n)
 \equiv
 \varepsilon_\lambda^\mu\varepsilon_\lambda^{\nu *}
 =\frac12\left[
 \mathsf P_\perp^{\mu\nu}(k;n)
 -\ii\lambda\mathsf C^{\mu\nu}(k;n)
 \right].}
 \end{equation*}
立刻有
\begin{equation*}
 \mathsf P_{+}^{\mu\nu}+\mathsf P_{-}^{\mu\nu}
 =\mathsf P_\perp^{\mu\nu}.
 \end{equation*}
因此未偏振光子求和正是两个 Wigner 螺旋度扇区的完备和。
\end{derivation}

\begin{derivation}[从\eqnrefs{eq:qed-S0-SI-dp7}到\eqnrefs{eq:qed-vertex-momentum-conservation-dp68}]
只看一次相互作用插入，并从三个自由场各取一个确定的平面波 Fourier 分量。把与 $x$ 无关的展开系数分别记为 $C_p,C_{p'},C_q$，则可严格写成
\begin{align*}
 \psi_p(x)&=C_pu(p)\ee^{-\ii p\cdot x/\hbar},\\
 \bar\psi_{p'}(x)&=C_{p'}^*\bar u(p')\ee^{+\ii p'\cdot x/\hbar},\\
 A_{q\mu}(x)&=C_q\varepsilon_\mu(q)\ee^{-\ii q\cdot x/\hbar},
\end{align*}
其中光子四动量 $q$ 取为流入顶角。三者代入\eqnrefs{eq:qed-S0-SI-dp7} 后，与 $x$ 有关的因子唯一合并成

\begin{equation*}
 \exp\!\left[
 \frac{\ii}{\hbar}(p'-p-q)\cdot x
 \right].
\end{equation*}
因此顶角时空积分是
\begin{align*}
 \int\dd^4 x\,
 \ee^{\frac{\ii}{\hbar}(p'-p-q)\cdot x}
 &=(2\pi\hbar)^4
 \delta^{(4)}(p'-p-q),\\
 &\Longrightarrow\qquad p+q=p'.
 \end{align*}
第一行只是全书 Fourier 约定下的 $\delta$-函数恒等式；第二行是它的支撑条件。也就是说，四动量守恒直接来自相互作用在时空中的平移不变性与局域积分。

内部线的四动量可以不满足 $q^2=m^2c^2$，因为它不是 LSZ 外态；但它连接的每个顶角仍各自携带一个这样的 \(\delta\) 函数。所以 ``虚粒子可以暂时违反能量守恒'' 并不是 QED 的数学陈述。真正被放宽的是内部动量的在壳约束，而不是顶角守恒律。
\end{derivation}

\begin{derivation}[从\eqnrefs{eq:qed-gell-mann-low-dp7}到\eqnrefs{eq:si-reduced-vertices-r10}]
先定义由自由 Dirac 与规范固定 Maxwell 二次作用量产生的归一化高斯平均

\begin{equation*}
 \langle\mathcal O\rangle_0
 \equiv
 \frac{\displaystyle
 \int\mathcal D\bar\psi\,\mathcal D\psi\,\mathcal DA\,
 \mathcal O\,\ee^{\ii S_0/\hbar}}
 {\displaystyle
 \int\mathcal D\bar\psi\,\mathcal D\psi\,\mathcal DA\,
 \ee^{\ii S_0/\hbar}}.
 \end{equation*}
自由高斯已在前面分别给出 Dirac 传播子 $S_F$ 与光子传播子 $D_F^{\mu\nu}$。相互作用理论的时序关联函数可以写成
\begin{equation*}
 {
 \langle\Tord\mathcal O\rangle
 =\frac{
 \left\langle\Tord\left\{\mathcal O
 \exp\!\left(\ii S_I/\hbar\right)\right\}\right\rangle_0}
 {\left\langle\Tord\exp\!\left(\ii S_I/\hbar\right)\right\rangle_0}.}
 \end{equation*}
分母的作用是消去与外部插入完全不连接的真空 bubbles；这正是连通 correlation 函数与真空归一化的关系。

为了看见基本顶角，不直接背诵 ``$-\ii e\gamma^\mu$''，而计算一个带两条费米子外插入和一条光子外插入的三点函数：
\begin{equation*}
 G_3^\mu(x,z;y)
 \equiv
 \left\langle
 \Tord\{\psi(x)A^\mu(z)\bar\psi(y)\}
 \right\rangle .
 \end{equation*}
把\eqnrefs{eq:qed-gell-mann-low-dp7} 的分子展开到 $q_{\mathrm e}$ 一阶，连通部分为
\begin{align*}
 G_{3,c}^{\mu,(1)}(x,z;y)
 ={}&-\frac{\ii q_{\mathrm e}}{\hbar}
 \int\dd^4 w\,
 \Big\langle\Tord\big\{
 \psi(x)A^\mu(z)\bar\psi(y)
 \bar\psi(w)\gamma^\nu\psi(w)A_\nu(w)
 \big\}\Big\rangle_{0,c}.
 \end{align*}
自由高斯中要让全部外部插入与相互作用点 $w$ 连通，只有一类收缩：
\begin{align*}
 \bigl[\psi(x)\,\bar\psi(w)\bigr]_{\rm contr}&\longrightarrow S_F(x-w),\\
 \bigl[\psi(w)\,\bar\psi(y)\bigr]_{\rm contr}&\longrightarrow S_F(w-y),\\
 \bigl[A^\mu(z)A_\nu(w)\bigr]_{\rm contr}&\longrightarrow D_F^{\mu}{}_{\nu}(z-w),
\end{align*}
下标 $\rm contr$ 表示选取这两个自由场的 Wick 收缩。
因此
\begin{equation*}
 {
 G_{3,c}^{\mu,(1)}(x,z;y)
 =-\frac{\ii q_{\mathrm e}}{\hbar}
 \int\dd^4 w\,
 S_F(x-w)\gamma^\nu S_F(w-y)
 D_F^{\mu}{}_{\nu}(z-w).}
 \end{equation*}
这条式子已经把图的三个组成部分全部暴露出来：$w$ 是顶角时空积分；三条自由收缩是三条传播子；唯一不属于自由传播的矩阵结构是 $\gamma^\nu$，它来自相互作用密度中的 Dirac 流。

Fourier 变换后，$w$ 积分产生四动量守恒 \(\delta\) 函数。再由 LSZ 去掉三条外部传播子极点，物理 SI 顶角因子就是
\begin{equation*}
 -\frac{\ii q_{\mathrm e}}{\hbar}\gamma^\nu,
 \end{equation*}
而把统一的 SI 场归一化常数提出后，约化规则变成前面定义的
\begin{equation*}
 {-\ii g_e\gamma^\nu,\qquad
 g_e=\frac{q_{\mathrm e}}{\sqrt{\varepsilon_0\hbar c}}.}
 \end{equation*}
所以 QED 顶角不是图论公设，而是 Maxwell--Dirac 作用量的一次微扰插入经过 Wick 收缩与 LSZ 后留下的截肢局域核。
\end{derivation}



\begin{derivation}[从\eqnrefs{eq:qed-loop-euler-dp7}到\eqnrefs{eq:qed-vertex-count-order-dp76}]

只含 $\bar\psi\psi A$ 基本顶角的连通 QED 图中，每个顶角提供两个费米子 half-edge 和一个光子 half-edge。若内部费米子线、内部光子线分别有 $I_f,I_\gamma$ 条，而外部费米子腿、光子腿分别有 $E_f,E_\gamma$ 条，则 half-edge 计数给
\begin{align*}
2V&=2I_f+E_f,\\
V&=2I_\gamma+E_\gamma.
\end{align*}
因此
\begin{equation*}
I_f=V-\frac{E_f}{2},
\qquad
I_\gamma=\frac{V-E_\gamma}{2}.
\end{equation*}
代入正文的连通图圈数关系~\eqnrefs{eq:qed-loop-euler-dp7}，
\begin{align*}
L
&=I_f+I_\gamma-V+1,\\
&=V-\frac{E_f}{2}+\frac{V-E_\gamma}{2}-V+1,\\
&=\frac{V-E_f-E_\gamma}{2}+1.
\end{align*}
解出
\begin{equation*}
V=2L+E_f+E_\gamma-2.
\end{equation*}
每个基本顶角带一个 $g_{\rm em}$，并采用前文固定的耦合约定
\begin{equation*}
 g_{\rm em}^2=4\pi\alpha.
\end{equation*}
把给定连通图中除显式耦合因子外的自旋、动量与圈积分部分定义为约化系数 $\widetilde{\mathcal M}_{L,E_f,E_\gamma}$，则
\begin{align*}
 \mathcal M_{L,E_f,E_\gamma}
 &=g_{\rm em}^{V}\widetilde{\mathcal M}_{L,E_f,E_\gamma}\\
 &=(4\pi\alpha)^{L+(E_f+E_\gamma)/2-1}
 \widetilde{\mathcal M}_{L,E_f,E_\gamma}.
\end{align*}
因此该图的显式耦合阶正是正文\eqnrefs{eq:qed-vertex-count-order-dp76}。
\end{derivation}
